\section{Conclusions}
\label{sec:conclusions}

We introduced automated code clone \emph{generation} as the constructive inverse of clone detection, and we presented \sys, an agentic system that synthesizes a behaviorally faithful but structurally divergent clone of a given program. \sys decomposes the task across four cooperating LLM-driven agents---library substitution, glue synthesis, first-party logic mutation, and cross-cutting verification---and enforces correctness at two tiers: a local execute-and-repair loop that guards every transformation, and a global equivalence oracle that certifies the finished clone or localizes its divergence. The key insight is that generation, unlike detection, permits constructive verification: an LLM need not be correct in one shot, because a disciplined loop can execute, test, and repair its output until equivalence holds.

On software that estimates factory measurements and controls industrial infrastructure---a domain where behavioral fidelity is a hard constraint---\sys generates equivalence-confirmed clones for 16 of 18 Category~C and~D runs at a median of zero repair iterations, and reduces the share of accepted clones carrying serious behavioral defects from 54\% to 33\% relative to an unverified single-shot baseline. Two caveats belong with that number. The difference is in the predicted direction but this corpus is too small to establish it ($p=0.34$, Fisher's exact). And \sys accepts far fewer clones than the baseline, so it delivers slightly fewer sound ones overall: the gain is in the precision of its acceptance signal, not in throughput. A configuration that trades recall for precision is the right default in a domain where a plausible-looking but wrong clone is worse than no clone, but it is a trade rather than a dominance result.

The ablation sharpens the attribution. \sys with its verification tier disabled behaves almost identically to the single-shot baseline---55\% against 54\%---which indicates that the agent decomposition on its own contributes little, and that what separates the full system is the verification discipline rather than the division of labor among agents.

Three directions follow directly. First, we will strengthen the oracle from test- and trace-based checking toward partial formal equivalence on the safety-critical core of control logic, tightening the coverage-bounded guarantee discussed in Section~\ref{sec:results}; our evaluation found defects that every automated layer accepted, so the bound is not merely theoretical. Second, the residual failure mode is concentrated at the seam between agents---adapters generated correctly and then incompletely wired to their call sites---which is a program-analysis problem rather than a generation problem, and suggests pairing each agent with a static check on the contract it is expected to honor. Third, because the capacity to produce equivalent-but-divergent code is inherently dual-use---defensively valuable for diversification, yet potentially useful for evading similarity-based detection---we will pair every generated clone with the machine-checkable provenance report \sys already emits, so that each transformation is attributable and auditable. Recognizing and generating clones are two faces of the same problem; coupling them, as \sys does, turns program equivalence from a question we merely answer into a property we can deliberately engineer.
